\section*{الفصل \ref{ch:g12:proton-nmr} --- الرنين المغناطيسي النووي للبروتون}

\begin{solution}{exo:g12:proton-nmr:1}
الميثان: 1. الإيثان: 1 (فمجموعتا \ce{CH3} متماثلتان). البروبان: 2 (مجموعتا
\ce{CH3} متماثلتان، ومجموعة \ce{CH2} واحدة). الميثانول: 2 (\ce{CH3} و\ce{OH}).
الميثوكسي ميثان: 1.
\end{solution}

\begin{solution}{exo:g12:proton-nmr:2}
جاران: إشارة ثلاثية، 1:2:1. ثلاثة جيران: إشارة رباعية، 1:3:3:1.
لا جيران: إشارة أحادية.
\end{solution}

\begin{solution}{exo:g12:proton-nmr:3}
المجموع \qty{40}{mm} لثمانية \omterm{def:g9:inside-the-atom:proton}{بروتونات}: \qty{5}{mm} لكل \omterm{def:g9:inside-the-atom:proton}{بروتون}. فالإشارات
تمثل 3 و2 و3 \omterm{def:g9:inside-the-atom:proton}{بروتونات}.
\end{solution}

\begin{solution}{exo:g12:proton-nmr:4}
\omterm{def:g11:functional-groups:carbonyl}{الألدهيد}: من 9.7 إلى \qty{10.0}{ppm}. ومجموعة \ce{CH3} في سلسلة: من 0.7 إلى
\qty{1.3}{ppm}.
\end{solution}

\begin{solution}{exo:g12:proton-nmr:5}
بروتوناته الاثنا عشر المتكافئة تعطي خطًا واحدًا حادًا شديدًا، يسهل
وضعه عند الصفر. \omterm{def:g6:solutions-and-solubility:solute}{والمذيب} المُدَوْتَر لا يعطي إشارة بروتونية خاصة
به، كانت ستطغى لولا ذلك على إشارات المليغرامات القليلة من العيّنة.
\end{solution}

\begin{solution}{exo:g12:proton-nmr:6}
إشارتان. \ce{CH2} المجاورة للكلور: بروتونان، إشارة رباعية (3
جيران)، في المجال 2.5--\qty{4.0}{ppm}. \ce{CH3}: 3 \omterm{def:g9:inside-the-atom:proton}{بروتونات}، إشارة
ثلاثية (جاران)، قرب المجال المعتاد لمجموعة \ce{CH3}، وأعلى منه قليلًا
لأن الكلور على بعد ذرتين.
\end{solution}

\begin{solution}{exo:g12:proton-nmr:7}
ثلاث إشارات: مجموعتا \ce{CH3}، 6 \omterm{def:g12:proton-nmr:equivalent}{بروتونات متكافئة}، إشارة ثنائية (جار
واحد، هو \ce{CH})؛ و\ce{CH}، \omterm{def:g9:inside-the-atom:proton}{بروتون} واحد، تقسمه 6 جيران
إلى 7 خطوط، في مجال \ce{H-C-O} (3.3--\qty{4.5}{ppm})؛
و\ce{OH}، \omterm{def:g9:inside-the-atom:proton}{بروتون} واحد، إشارة أحادية.
\end{solution}

\begin{solution}{exo:g12:proton-nmr:8}
(a) البروبانون؛ (b) إيثانوات الإيثيل؛ (c) حمض البروبانويك.
\end{solution}

\begin{solution}{exo:g12:proton-nmr:9}
\ce{CH2} مرتبطة \omterm{def:g7:atoms-and-molecules:atom}{بذرة} أكسجين \omterm{def:g11:functional-groups:carboxylic-acid}{الإستر}، وهي
\omterm{def:g11:polarity-and-cohesion:electronegativity}{كهرسلبية} فتنقص حجب بروتوناتها: مجال \ce{H-C-O}،
3.3--\qty{4.5}{ppm}. أما \ce{CH3} فأبعد \omterm{def:g7:atoms-and-molecules:atom}{بذرة} كربون واحدة وتبقى في
مجال السلاسل.
\end{solution}

\begin{solution}{exo:g12:proton-nmr:10}
البروبانون: مجموعة \ce{C=O} (\omterm{def:g11:functional-groups:carbonyl}{كيتون}، \qty{1715}{cm^{-1}}) وستة \omterm{def:g12:proton-nmr:equivalent}{بروتونات
متكافئة}، إشارة أحادية واحدة. أما البروبانال فيعطي ثلاث إشارات، إحداها،
\omterm{def:g9:inside-the-atom:proton}{بروتون} \omterm{def:g11:functional-groups:carbonyl}{الألدهيد} في \ce{CH3-CH2-CHO}، نحو 9.7--\qty{10.0}{ppm}.
\end{solution}

\begin{solution}{exo:g12:proton-nmr:11}
\omterm{def:g9:inside-the-atom:proton}{بروتون} \ce{OH} يتبادل بسرعة بين \omterm{def:g7:atoms-and-molecules:molecule}{الجزيئات}: فيعطي إشارة
أحادية ولا يقسم إشارات جيرانه. لذلك لا تقسم إشارةَ \ce{CH2}
إلا بروتوناتُ \ce{CH3} الثلاثة: إشارة رباعية.
\end{solution}

\begin{solution}{exo:g12:proton-nmr:12}
حمض البوتانويك \ce{CH3-CH2-CH2-COOH}: 4 إشارات، ثلاثية (\ce{CH3})،
وإشارة من 6 خطوط (\ce{CH2} الوسطى، 5 جيران)، وثلاثية
(\ce{CH2-C=O})، وأحادية نحو 11--\qty{12}{ppm}. إيثانوات الإيثيل: أحادية
ورباعية وثلاثية. بروبانوات الميثيل \ce{CH3-CH2-CO-O-CH3}:
ثلاثية، ورباعية (\ce{CH2-C=O}، 2.0--\qty{2.4}{ppm})، وأحادية
(\ce{O-CH3}). ميثانوات البروبيل \ce{H-CO-O-CH2-CH2-CH3}: 4 إشارات،
أحادية بعيدة إلى اليسار (H على كربون \ce{C=O})، وثلاثية
(\ce{O-CH2})، وإشارة من 6 خطوط، وثلاثية (\ce{CH3}).
\end{solution}

\begin{solution}{exo:g12:proton-nmr:13}
إشارة \ce{OH}: \omterm{def:g9:inside-the-atom:proton}{فبروتونات} \ce{O-H} تتبادل مع دوتيريوم
الماء الثقيل، و\ce{O-D} لا تعطي إشارة بروتونية. فالإشارة التي
تختفي بالرج مع \ce{D2O} تعود إلى \omterm{def:g9:inside-the-atom:proton}{بروتون} على O أو N.
\end{solution}

\begin{solution}{exo:g12:proton-nmr:14}
يضيف الإيثانول إشارة رباعية نحو \qty{3.69}{ppm}، وثلاثية نحو
\qty{1.23}{ppm} وإشارة \ce{OH} أحادية. الإشارة الأحادية \omterm{def:g11:functional-groups:carboxylic-acid}{للإستر}: \qty{30}{mm}
لثلاثة \omterm{def:g9:inside-the-atom:proton}{بروتونات}، أي \qty{10}{mm} لكل \omterm{def:g9:inside-the-atom:proton}{بروتون} من \omterm{def:g11:functional-groups:carboxylic-acid}{الإستر}. والإشارة الثلاثية للإيثانول:
\qty{3}{mm} لثلاثة \omterm{def:g9:inside-the-atom:proton}{بروتونات}، أي \qty{1}{mm} لكل \omterm{def:g9:inside-the-atom:proton}{بروتون} من الإيثانول. ولأن كل
\omterm{def:g7:atoms-and-molecules:molecule}{جزيء} يسهم بمجموعة \ce{CH3} واحدة في هاتين الإشارتين، فالكميتان بنسبة
$1/10$: نحو \omterm{def:g7:atoms-and-molecules:molecule}{جزيء} إيثانول واحد لكل عشرة \omterm{def:g7:atoms-and-molecules:molecule}{جزيئات} من \omterm{def:g11:functional-groups:carboxylic-acid}{الإستر}.
\end{solution}

\begin{solution}{exo:g12:proton-nmr:15}
أربع إشارات: مجموعتا \ce{CH3} (6 \omterm{def:g9:inside-the-atom:proton}{بروتونات}) إشارة ثنائية؛ و\ce{CH} (\omterm{def:g9:inside-the-atom:proton}{بروتون}
واحد) تقسمه 6 + 2 = 8 جيران إلى 9 خطوط؛ و\ce{CH2} (بروتونان)
المجاورة \omterm{def:g7:atoms-and-molecules:atom}{لذرة} \ce{O}، إشارة ثنائية، في مجال \ce{H-C-O}؛
و\ce{OH} (\omterm{def:g9:inside-the-atom:proton}{بروتون} واحد) إشارة أحادية. وإشارة \ce{CH} لها أكبر عدد من الخطوط.
\end{solution}

\begin{solution}{pb:g12:proton-nmr:1}
\textbf{1.} \omterm{def:g11:organic-skeletons:alkane}{الألكان} ذو أربع \omterm{def:g7:atoms-and-molecules:atom}{ذرات} كربون هو \ce{C4H10}؛ ومجموعة \ce{C=O} واحدة
تنقص ذرتي هيدروجين، وذرتا الأكسجين لا تضيفان شيئًا: \ce{C4H8O2}، وهي
صيغة الأحماض \omterm{def:g11:functional-groups:carboxylic-acid}{والإسترات} ذات الأربع \omterm{def:g7:atoms-and-molecules:atom}{ذرات} كربون.

\textbf{2.} حمض البوتانويك \ce{CH3-CH2-CH2-COOH}؛ إيثانوات الإيثيل
\ce{CH3-CO-O-CH2-CH3}؛ بروبانوات الميثيل \ce{CH3-CH2-CO-O-CH3}؛ ميثانوات
البروبيل \ce{H-CO-O-CH2-CH2-CH3}.

\textbf{3.} تشترك \omterm{def:g11:functional-groups:carboxylic-acid}{الإسترات} الثلاثة في مجموعة \omterm{def:g11:functional-groups:carboxylic-acid}{الإستر}؛ وللأربعة كلها
مجموعة \ce{C=O}.

\textbf{4.} مجموعة \ce{C=O}، في مجال \omterm{def:g11:functional-groups:carboxylic-acid}{الإسترات} (نحو \qty{1735}{cm^{-1}}).

\textbf{5.} حزمة \ce{O-H} العريضة جدًا من 2500 إلى
\qty{3300}{cm^{-1}}: وهي غائبة، فليس المركب حمض البوتانويك.

\textbf{6.} لا: فللإسترات الثلاثة المجموعات نفسها.

\textbf{7.} ثلاث.

\textbf{8.} مجموعة إيثيل \ce{CH3-CH2-}: \ce{CH2} تقسمها 3
\omterm{def:g9:inside-the-atom:proton}{بروتونات} (رباعية)، و\ce{CH3} يقسمها بروتونان (ثلاثية).

\textbf{9.} ليس لبروتوناتها \omterm{def:g9:inside-the-atom:proton}{بروتونات} على \omterm{def:g7:atoms-and-molecules:atom}{الذرات} المجاورة: فهي
\ce{CH3} مرتبطة بكربون \ce{C=O} أو بأكسجين \omterm{def:g11:functional-groups:carboxylic-acid}{الإستر}.

\textbf{10.} \ce{CH2} مرتبطة بالأكسجين (مجال \ce{H-C-O}).

\textbf{11.} إشارة ثلاثية (\ce{CH3})، وإشارة رباعية نحو 2.0--\qty{2.4}{ppm} تعود إلى
\ce{CH2} المجاورة لمجموعة \ce{C=O}، وإشارة أحادية في مجال \ce{H-C-O} تعود إلى
\ce{O-CH3}. فالرباعية ستكون أدنى بكثير من \qty{4.12}{ppm} والأحادية
أعلى بكثير من \qty{2.04}{ppm}: فليس هو المجهول.

\textbf{12.} أربع إشارات (أحادية بعيدة إلى اليسار، وثلاثية، وإشارة من 6
خطوط، وثلاثية): فليس هو المجهول، الذي له ثلاث.

\textbf{13.} إيثانوات الإيثيل، \ce{CH3-CO-O-CH2-CH3}: الأحادية عند
\qty{2.04}{ppm} = \ce{CH3-C=O}؛ والرباعية عند \qty{4.12}{ppm} = \ce{O-CH2}؛
والثلاثية عند \qty{1.26}{ppm} = \ce{CH3} الطرفية.

\textbf{14.} الأحادية 3، والرباعية 2، والثلاثية 3.

\textbf{15.} $72 / 8 = \qty{9}{mm}$ لكل \omterm{def:g9:inside-the-atom:proton}{بروتون}.

\textbf{16.} الأحادية $3 \times 9 = \qty{27}{mm}$؛ والثلاثية
\qty{27}{mm}.

\textbf{17.} $18 + 27 + 27 = \qty{72}{mm}$.

\textbf{18.} نعم: \omterm{def:g11:functional-groups:carboxylic-acid}{فالإسترات} الصغيرة معروفة بروائحها الفاكهية.

\textbf{19.} درجة الإشارة الرباعية \qty{18}{mm} من أصل \qty{72}{mm}.
\end{solution}
